\documentclass[11pt,a4paper]{article}

% ======================================================================
%  Erlangen 2026-05-28 AutoPos OptiTrack Validation -- LONG SUMMARY
%  Structure assumed:  main.tex  +  figs/*.png
%  All figures are referenced as figs/<name>.png with the original names.
% ======================================================================

\usepackage[utf8]{inputenc}
\usepackage[T1]{fontenc}
\usepackage{lmodern}
\usepackage[margin=2.4cm]{geometry}
\usepackage{graphicx}
\usepackage{booktabs}
\usepackage{amsmath,amssymb}
\usepackage{siunitx}
\usepackage{longtable}
\usepackage{array}
\usepackage{xcolor}
\usepackage{colortbl}
\usepackage{caption}
\usepackage{subcaption}
\usepackage{enumitem}
\usepackage{fancyhdr}
\usepackage{microtype}
\usepackage[hidelinks,colorlinks=true,linkcolor=blue!50!black,citecolor=blue!50!black,urlcolor=blue!50!black]{hyperref}

\graphicspath{{figs/}}

\sisetup{detect-weight=true,detect-family=true}
\newcommand{\mm}{\,\si{\milli\meter}}
\newcommand{\todo}[1]{{\color{red}\textbf{[#1]}}}
\newcommand{\note}[1]{{\color{blue!50!black}\emph{#1}}}

\pagestyle{fancy}
\fancyhf{}
\rhead{Erlangen AutoPos Validation --- Long Summary}
\lhead{2026-05-28 Official}
\cfoot{\thepage}

\title{\textbf{AutoPos Pure-UWB 3D Anchor Self-Calibration:\\
Absolute Validation Against OptiTrack\\[4pt]
\large Erlangen Official Experiment, 2026-05-28 --- Detailed Technical Report (Long Summary)}}
\author{Zekai Xiao}
\date{Generated 2026-05-30}

\begin{document}
\maketitle
\thispagestyle{fancy}

\begin{abstract}
\noindent
This report documents the absolute validation of the AutoPos pure-UWB anchor
self-calibration system against an independent OptiTrack optical reference, using the
official Erlangen dataset captured on 2026-05-28. AutoPos recovers the full 3D layout of
eight UWB anchors (DWM1001C; DW1000 + nRF52832) from \emph{inter-anchor ranging only},
fixing anchor~A at the origin and jointly estimating per-anchor antenna delays. We report:
(i) anchor-layout absolute accuracy after reflection-allowed rigid alignment
($\approx\!105\mm$ RMS); (ii) static-tag absolute accuracy under the AutoPos layout
($\approx\!77\mm$ median 3D, production output), with a strongly anisotropic error
structure (horizontal P95 $\approx\!58$--$80\mm$ versus vertical P95 $\approx\!260\mm$);
(iii) a surveyed-anchor control isolating self-calibration cost from the UWB ranging floor;
(iv) range-only VDOP geometry; (v) MC5000 and stratified keep-$k$ dropout robustness;
(vi) per-anchor diagnostics (pair residuals, temporal drift, health scorecard); and
(vii) an error-structure synthesis attributing the dominant tag error to radial,
scale-coupled layout error rather than to one weak anchor or isotropic noise. A dedicated
section is reserved for rotational (roto) absolute validation, which is pending external
OptiTrack processing. The honest headline is that AutoPos delivers credible decimetre-level
3D anchor calibration and approximately $6$--$8\,\si{\centi\meter}$ static absolute capture
with a fully auditable error budget, while not claiming to be a general centimetre-level
optical-mocap replacement.
\end{abstract}

\tableofcontents
\clearpage

% ======================================================================
\section{Introduction and Scope}
% ======================================================================

\subsection{Purpose of this document}
This is the \emph{long} version of the Erlangen validation report. It is intentionally
exhaustive: it records every analysis performed, the conventions and pitfalls encountered,
the numerical results, and the honest interpretation of each result, including null results
and corrections of earlier expectations. A condensed short summary for the supervisor will
be derived from this document; nothing here is written to be persuasive at the expense of
accuracy.

\subsection{System under test}
AutoPos is a pure-UWB anchor self-calibration pipeline. Eight anchors, labelled A--H, are
built on the DWM1001C module (Decawave DW1000 transceiver plus Nordic nRF52832). The
anchors exchange broadcast single-sided two-way ranging (SS-TWR) measurements. From the
resulting inter-anchor range set \emph{alone} --- with no external position reference at
deployment time --- the solver estimates the 3D coordinates of all anchors, fixing
A~$=(0,0,0)$ as the gauge origin and estimating per-anchor antenna delays differentially
(A's delay fixed at zero). The deployment uses a two-layer arrangement (a lower set
A--D and an upper set E--H) approximating a rectangular box; the layers are separated by
roughly \SI{1.4}{\meter} in the vertical direction.

\subsection{Why this validation matters}
Internal repeatability tells us whether repeated estimates cluster tightly in the solver's
own frame; it does not tell us whether the resulting coordinates are \emph{correct} in the
real world. The Erlangen campaign provides an independent optical ground truth (OptiTrack,
sub-millimetre marker noise) so that AutoPos's self-reported ``repeatability'' can be
converted into \emph{absolute} accuracy, and so that the error can be decomposed into its
physical sources.

\subsection{What is and is not claimed}
\begin{itemize}[leftmargin=1.4em]
  \item We claim a validated, auditable 3D anchor layout at the $\sim\!105\mm$ RMS level
        and static-tag absolute accuracy at the $\sim\!62$--$77\mm$ median level inside the
        working volume.
  \item We do \emph{not} claim a general centimetre-level optical-mocap replacement, nor a
        validated absolute rotational accuracy (see Section~\ref{sec:roto}).
\end{itemize}

% ======================================================================
\section{Coordinate Frames and Validation Conventions}
% ======================================================================
\label{sec:conventions}

Two coordinate frames are used throughout and must never be compared component-wise without
an explicit rigid alignment.

\paragraph{AutoPos solver frame.}
\[
  x,\,y \;=\; \text{horizontal plane}, \qquad
  z \;=\; \text{vertical axis (upper layer is negative $z$)}, \qquad
  \text{display height} = -z.
\]

\paragraph{OptiTrack TRC frame.}
\[
  Y \;=\; \text{vertical}, \qquad X,\,Z \;=\; \text{horizontal}.
\]
This was verified empirically, not assumed: the upper and lower anchor layers separate by
$\approx\!1423\mm$ along OptiTrack $Y$ and by only $5$--$8\mm$ along $X$ and $Z$. Hence
$Y$ is unambiguously the vertical axis in the optical frame. A common implementation pitfall
is to reuse the AutoPos ``$z$ is vertical'' convention on OptiTrack data; doing so would
silently take a horizontal axis as vertical and corrupt every per-axis result. A dedicated
audit confirmed that all per-axis splits in the OptiTrack frame use $Y$ (index 1) as
vertical.

\paragraph{Consequences for alignment.}
\begin{itemize}[leftmargin=1.4em]
  \item Layout comparison requires a full rigid alignment between the AutoPos and OptiTrack
        frames.
  \item The alignment must \textbf{allow reflection}. Inter-anchor distance calibration
        cannot determine handedness (the distance matrix is reflection-invariant), and the
        AutoPos layout converges to the mirror branch relative to OptiTrack. Proper-rotation
        Kabsch yields nonsensical $\approx\!1.4$--$1.5\,\si{\meter}$ RMS; reflection-allowed
        Kabsch yields $\approx\!105\mm$.
  \item The official accuracy figure is the rigid, \emph{no-scale} RMS. A similarity
        (scale-fitted) RMS is reported only as a diagnostic; it uses OptiTrack truth to
        remove scale and therefore must not be quoted as an accuracy claim.
  \item Static-tag absolute accuracy uses a transform locked \emph{from the anchors only};
        the tag ground truth is never used to fit the transform that later measures tag
        error (this would be circular).
\end{itemize}

% ======================================================================
\section{Dataset and Solver Matrix}
% ======================================================================

\paragraph{Anchor-layout solver versions.} \texttt{v1-old}, \texttt{v2}, \texttt{v3-lite},
\texttt{v3-full}, \texttt{v4-io}. \texttt{v4-io} is the current production solver.

\paragraph{Tag solver family.} \texttt{T1}, \texttt{T2}, \texttt{T3}, \texttt{T4}.
\texttt{T4} is the robustness-oriented deployment solver; \texttt{T3} gives the best static
median on this dataset.

\paragraph{Static-tag test grid.} Twenty-four static positions (ID01--ID24), spanning
edge and centre placements, three heights (low/mid/high), and four facing groups
(ABEF/BCGF/CDHG/ADHE), so that all faces of the array are exercised rather than only the
positions nearest a few anchors. The tag fixture (\texttt{BSF66F}) carries five glued
optical balls plus the virtual \texttt{Icenter}/\texttt{Iantenna} markers.

% ======================================================================
\section{Methods}
% ======================================================================
\label{sec:methods}

This section describes the processing chain used to produce the validation results. It is
written as an experimental method, not as a general UWB theory chapter. The important point
is which data were used, which degrees of freedom were estimated, which quantities were held
out for validation, and which operations were deliberately avoided to prevent circularity.

\subsection{End-to-end analysis pipeline}
The full chain has four data layers:
\begin{enumerate}[leftmargin=1.6em]
  \item \textbf{Raw UWB ranging.} Anchor-to-anchor sweep data are used for AutoPos layout
        self-calibration. Static tag-to-anchor captures are used for tag localization
        validation. The layout solver does not consume OptiTrack positions.
  \item \textbf{OptiTrack reference.} Optical marker exports provide anchor antenna truth
        and static tag truth. These are used only for validation, frame alignment, and the
        explicitly labelled surveyed-anchor control.
  \item \textbf{AutoPos layout outputs.} Each anchor-layout solver variant writes a
        \texttt{layout.json} with anchor coordinates and per-anchor delay terms
        \texttt{d\_anchor\_mm}. These layouts feed the tag solvers.
  \item \textbf{Report-side post-processing.} Python scripts under
        \texttt{official\_extra\_analysis/scripts} compute alignment metrics, static-tag
        error, CDF/outlier metrics, DOP maps, MC summaries, drift diagnostics, and
        additional error-structure diagnostics. Each script logs provenance in
        \texttt{run\_meta.json}.
\end{enumerate}

\subsection{Traditional SS-TWR and broadcast SS-TWR data roles}
\label{sec:sstwrmethod}
The system contains two relevant single-sided two-way ranging (SS-TWR) usage modes. The
report does not re-derive SS-TWR timing theory; it only states how each mode is used in this
dataset.

\paragraph{Traditional point-to-point SS-TWR.}
Traditional SS-TWR is the direct pairwise ranging mode: one initiator exchanges messages
with one responder and obtains one range observation for that pair. It is conceptually the
cleanest way to think about a range measurement and is useful for calibration stations,
simple sanity checks, and tag-to-anchor captures. In the static tag captures, each frame is
ultimately consumed as a set of tag-to-anchor ranges
\[
  \{(\text{anchor }i,\; \rho_i)\}_{i\in \mathcal{A}_t},
\]
where $\mathcal{A}_t$ is the set of anchors valid in that frame. The tag solver later
corrects each measured range using the selected layout's anchor delay and an optional tag
delay term.

\paragraph{Broadcast SS-TWR.}
The anchor self-calibration data are generated from a broadcast-style SS-TWR sweep: anchors
participate in scheduled exchanges so that one capture pass yields a dense matrix of
inter-anchor ranges rather than a single pair at a time. The report treats the final
artifact as a set of robust pair medians
\[
  \hat{\rho}_{ij}, \qquad 1\le i < j \le 8,
\]
and, for diagnostic plots, the raw directional measurements are also retained to estimate
asymmetry and pair stability. Broadcast SS-TWR is therefore a data-acquisition efficiency
choice: it creates the dense inter-anchor observation graph needed by AutoPos without
requiring manual pair-by-pair measurement. All downstream layout solvers see the resulting
pairwise range table, not the packet schedule itself.

\paragraph{Why both are mentioned.}
The layout problem and the tag problem are intentionally separated. Broadcast SS-TWR
supplies the anchor-to-anchor graph for self-calibration; traditional/static tag ranging
supplies tag-to-anchor frames for validation. Mixing these roles would create a circular
argument: static tag truth must validate the layout and tag solver, not train or align it.

\subsection{Anchor-layout model and gauge choices}
The anchor layout is estimated from inter-anchor ranges. For anchors $i,j$ with positions
$a_i,a_j$ and endpoint delay terms $d_i,d_j$, the residual used by delay-aware variants can
be summarized as
\[
  r_{ij} = \lVert a_i-a_j\rVert + d_i + d_j - \hat{\rho}_{ij}.
\]
The coordinate gauge is fixed internally by setting anchor A at the origin and choosing a
consistent orientation convention in the solver frame. The absolute pose, handedness, and
scale relative to OptiTrack are not known from UWB distances alone and are evaluated later
against the optical reference.

The \texttt{d\_anchor\_mm} values in \texttt{layout.json} are endpoint delay terms used by
the tag solver. They must be read as part of the joint layout solution. The later surveyed
baseline proves that the \texttt{v4-io} delay vector is not a portable physical calibration:
it is coupled to the self-calibrated geometry.

\subsection{Anchor-layout solver variants V1 to V4-io}
\label{sec:vmethod}
The V-series names refer only to anchor-layout solvers. They are separate from the T-series
tag solvers.

\begin{description}[leftmargin=2.6em,style=nextline]
  \item[\texttt{v1-old}.]
  Archived baseline using classical distance-geometry/MDS-style layout recovery with zero
  per-anchor delay. It is simple and important as a historical baseline. In the final
  absolute validation it remains surprisingly competitive: \texttt{v1-old} all8 rigid RMS
  is $106.9\mm$, only $\approx2\mm$ worse than \texttt{v4-io}.

  \item[\texttt{v2}.]
  Improved zero-delay layout solver variant. It reduces internal AutoPos pair residuals
  strongly relative to \texttt{v1-old} in the field-check tables, but after OptiTrack
  absolute alignment it is not the best layout. This is a useful reminder that internal
  residual quality and external absolute accuracy are related but not identical.

  \item[\texttt{v3-lite}.]
  Lightweight successor to \texttt{v2}, also zero-delay in the final exported layout. It
  behaves very similarly to \texttt{v2}: good internal pair residuals, but absolute
  OptiTrack RMS around $135\mm$ rather than the \texttt{v4-io} $\approx105\mm$.

  \item[\texttt{v3-full}.]
  More flexible delay-aware variant. It estimates nonzero per-anchor delay terms, including
  a wider delay span than \texttt{v4-io}. In this dataset the extra freedom does not improve
  the final absolute layout: all8 rigid RMS is $144.4\mm$. The interpretation is that
  additional freedom can absorb range inconsistencies in ways that look good internally but
  hurt the external geometry.

  \item[\texttt{v4-io}.]
  Current production layout solver. It is delay-aware, uses bounded endpoint delays
  (A fixed at $0\mm$, other anchors in this run between roughly $27$ and $60\mm$), and uses
  soft physical/two-layer priors appropriate for the deployment geometry. It is the best
  absolute layout in this dataset, but only marginally better than \texttt{v1-old};
  therefore it is reported as the preferred current solver, not as a dramatic standalone
  breakthrough.
\end{description}

\subsection{Tag positioning model and T1--T4 solver family}
\label{sec:tmethod}
The T-series names refer to tag positioning, not anchor layout. For each UWB frame, the
solver estimates the tag position $p$ from valid tag-to-anchor ranges:
\[
  r_i(p) = \lVert p-a_i\rVert + d_i + d_{\text{tag}} - \rho_i,
\]
where $a_i$ and $d_i$ come from the selected \texttt{layout.json}, and the default tag delay
is $d_{\text{tag}}=0$ unless a controlled baseline explicitly sets it.

\begin{description}[leftmargin=2.6em,style=nextline]
  \item[\texttt{T1}: robust WLS multilateration.]
  Per-frame 3D multilateration with Gauss--Newton iterations, anchor sigma weighting, and
  Huber robust residual weighting. It requires at least four valid anchors. The previous
  frame may be used as an initialization but is not a motion prior.

  \item[\texttt{T2}: quality-aware robust WLS.]
  Adds current quality weighting and a short quality exponential moving average. It does
  not hard-reject anchors; it changes their effective sigma.

  \item[\texttt{T3}: dynamic-stable robust WLS.]
  Adds persistent-residual soft downweighting and a weak previous-position prior. The goal
  is stable dynamic tracking under low redundancy, not single-frame residual minimization.
  It deliberately avoids aggressive leave-one-out rejection because prior experiments showed
  that hard rejection can reduce residual RMS while damaging trajectory consistency.

  \item[\texttt{T4}: adaptive redundancy policy.]
  Current deployment-oriented policy. With all eight anchors valid, T4 uses the memory-free
  T1 path to avoid carrying temporal bias. With fewer than eight anchors, it switches to
  T3-style stabilization and the weak previous-position prior. This matches the MC keep-$k$
  evidence: full-anchor frames have enough geometry, while keep-4/5/6 frames benefit from
  temporal stabilization.
\end{description}

The raw replay matrix evaluates all combinations
\[
  5\;\text{layout solvers}\times4\;\text{tag solvers}\times\{\text{all8},\text{noG}\}.
\]
The production-output headline is not automatically replaced by the best replay result:
\texttt{v4-io/T3/all8} is labelled best-case, while the representative production output
remains \texttt{v4-io/all8}.

\subsection{OptiTrack marker processing and corrected tag truth}
Anchor truth is obtained from the OptiTrack marker-derived antenna positions. Static tag
truth uses the virtual \texttt{Iantenna} marker. Because ID01 and ID05 have mislabeled
\texttt{I1..I5} optical balls, their tag truth is reconstructed as follows:
\begin{enumerate}[leftmargin=1.6em]
  \item take per-marker medians over frames for the five balls and \texttt{Iantenna};
  \item compare the ten inter-ball distances to the clean 22-capture consensus;
  \item apply the deterministic relabel permutations for ID01 and ID05;
  \item learn the physical ball-local antenna offset from clean captures;
  \item project the corrected antenna point back into each affected capture.
\end{enumerate}
All static-tag validation uses this corrected truth path.

\subsection{Alignment and metric computation}
Anchor layout accuracy uses reflection-allowed rigid Kabsch with no scale. Static-tag
accuracy uses the same anchor-locked transform and then compares solved tag coordinates to
corrected OptiTrack truth. The main 3D metric set includes mean, RMSE, P50, P90, P95, P99,
max, and cm-scale outlier rates at 50/80/100/200/300\,\si{\milli\meter}. Per-axis errors are
reported in the OptiTrack frame with $Y$ vertical; 2D horizontal error is secondary only.

\subsection{Surveyed-anchor and delay-control method}
The surveyed-anchor baseline replaces the AutoPos layout with OptiTrack-truth anchor
coordinates and solves the tag directly in the OptiTrack frame. Three delay treatments are
tested: zero delay, the AutoPos \texttt{v4-io} delay vector, and an OptiTrack-derived
inter-anchor endpoint delay fit. This baseline is explicitly a control: the delaycal row is
partly circular and should be read as an optimistic lower bound, while the AutoPos-delay row
tests whether the AutoPos delay vector is portable to a surveyed layout.

\subsection{Geometry, robustness, and diagnostic methods}
\begin{itemize}[leftmargin=1.4em]
  \item \textbf{VDOP/DOP.} The default DOP uses the range-only Jacobian row
        $[u_x,u_y,u_z]$, matching the solver. A range-bias Jacobian is appendix-only.
  \item \textbf{MC keep-$k$.} Random keep-$k$ repeats quantify robustness as anchors are
        forced unavailable, over solver/version/mode combinations.
  \item \textbf{Stratified keep-$k$.} Exhaustive fixed dropped-set replay separates
        lower-heavy, upper-heavy, balanced, and specific-anchor dropout effects.
  \item \textbf{Pair residual and raw asymmetry.} Inter-anchor residual heatmaps expose
        pair-specific bias; raw directional asymmetry checks whether a pair behaves
        differently depending on initiator/responder direction.
  \item \textbf{Temporal drift.} Static raw range links are regressed against time within
        each capture to measure drift slope and drift-over-capture.
  \item \textbf{Additional diagnostics.} Delay decomposition, tag error radial/tangential
        split, worst-point residual fingerprints, anchor health, height/edge/facing splits,
        and single-anchor criticality are computed from already-generated official outputs.
\end{itemize}

% ======================================================================
\section{Anchor-Layout Absolute Accuracy}
% ======================================================================
\label{sec:layout}

\subsection{Method}
AutoPos anchor coordinates are aligned to the OptiTrack antenna marker positions using
reflection-allowed Kabsch. Because of the mirror ambiguity (Section~\ref{sec:conventions}),
reflection must be permitted; the proper-rotation result is reported only as a chirality
sanity check.

\subsection{Headline results (\texttt{v4-io})}
\begin{table}[h]
\centering
\caption{Anchor-layout absolute accuracy, \texttt{v4-io}, reflection-allowed alignment.
The similarity (scale-fitted) RMS is a diagnostic only and is not an accuracy claim.}
\begin{tabular}{lrrrrrr}
\toprule
eval set & rigid RMS 3D & horiz.\ RMS & vert.\ RMS & shape RMS & sim.\ scale & sim.\ RMS \\
\midrule
all8 & 104.9 & 86.1 & 59.9 & 130.5 & 0.960 & 71.2 \\
noG  & 104.4 & 84.4 & 61.4 & 133.4 & 0.960 & 70.1 \\
\bottomrule
\end{tabular}
\end{table}

\subsection{Interpretation of the three layout numbers}
\begin{itemize}[leftmargin=1.4em]
  \item \textbf{Rigid RMS $=104.9\mm$} is the official absolute accuracy: best rigid fit
        (rotation + translation + reflection, no scale), then per-anchor residual RMS.
  \item \textbf{Shape RMS $=130.5\mm$} is alignment-free: it is the RMS difference of the 28
        inter-anchor pairwise distances between AutoPos and OptiTrack. It cannot be hidden by
        any choice of pose or handedness, so it is the most honest lower bound on the
        intrinsic shape distortion. It is larger than the rigid RMS precisely because it does
        not allow pose to absorb part of the error, and it includes the $\approx\!4\%$ scale
        bias (which inflates every pairwise distance).
  \item \textbf{Similarity RMS $=71.2\mm$} (scale 0.960) is what remains if a global scale is
        fitted using OptiTrack truth. This is diagnostic only: it quantifies how much of the
        error is scale, but quoting it as accuracy would mean using the ground truth to
        flatter the result.
\end{itemize}

\subsection{Solver ranking}
\begin{table}[h]
\centering
\caption{Solver ranking by reflection-allowed all8 rigid RMS.}
\begin{tabular}{lrrl}
\toprule
solver & rigid RMS 3D (\si{\milli\meter}) & shape RMS (\si{\milli\meter}) & note \\
\midrule
v4-io   & 104.9 & 130.5 & best, but only marginally \\
v1-old  & 106.9 & 137.6 & nearly tied with v4-io \\
v2      & 135.4 & 185.8 & worse \\
v3-lite & 135.4 & 185.7 & similar to v2 \\
v3-full & 144.4 & 149.2 & did not become best despite extra freedom \\
\bottomrule
\end{tabular}
\end{table}

\noindent
\textbf{Honest reading.} \texttt{v4-io} is the best layout solver but only by
$\approx\!2\mm$ over the simplest classical-MDS \texttt{v1-old}. This is too small to market
as a major algorithmic improvement on this dataset; the residual range structure and the
OptiTrack G-marker uncertainty (Sections~\ref{sec:anchorhealth}, \ref{sec:pair}) dominate
over the V1--V4 difference. The defensible statement is that \texttt{v4-io} is preferred
because it is best-or-near-best across the combined analysis, not because it is dramatically
more accurate in layout alone.

\begin{figure}[h]
\centering
\includegraphics[width=0.78\textwidth]{layout_opti_vs_autopos_3d.png}
\caption{AutoPos anchor layout versus OptiTrack ground truth after reflection-allowed rigid
alignment.}
\label{fig:layout3d}
\end{figure}

% ======================================================================
\section{Static-Tag Absolute Accuracy}
% ======================================================================
\label{sec:tag}

This is the system-output validation layer: it measures how much absolute tag-position error
results when solved tag positions are mapped into the OptiTrack frame using only the
anchor-derived transform.

\subsection{Frame-locking (why 8-anchor Procrustes, not a tag fit or a centroid)}
\label{sec:framelock}
The solved tag positions live in the AutoPos frame; the tag ground truth lives in the
OptiTrack frame. Mapping one to the other requires a transform $T$, and the choice of $T$
determines the error. $T$ must therefore be fixed by something \emph{other} than the tags.

\begin{table}[h]
\centering
\caption{Three ways to fix the frame transform. Only the anchor-locked method gives a
unique, non-circular tag error.}
\begin{tabular}{p{4.4cm}cccc}
\toprule
Method & locks transl. & locks rot. & locks mirror & tag error unique? \\
\midrule
(A) fit $T$ to the tag cloud         & yes & yes & yes & \textbf{no --- circular} \\
(B) align centroid / diagonal point  & yes & \textbf{no} & \textbf{no} & \textbf{no --- rotatable} \\
(C) 8-anchor Procrustes, then apply  & yes & yes & yes & \textbf{yes} \\
\bottomrule
\end{tabular}
\end{table}

\noindent
Method~(A) uses the error definition to minimise the error (circular; artificially small).
Method~(B), a single-point/centroid alignment, fixes only the three translation DOF and
leaves rotation and handedness free, so the tag cloud can be spun about the centre and every
orientation yields a different error --- it is \emph{less} determined than (A). Only
method~(C), where the transform is locked by all eight anchors (rigid, reflection-allowed,
no scale) and then applied to the tags, gives a unique and independent tag error. The
centroid is used only as an \emph{interpretation} origin (for error-versus-distance), never
to define the frame.

\subsection{Tag ground-truth correction: the ID01/ID05 marker-label swap}
\label{sec:swap}
The static-tag truth is the virtual marker \texttt{Iantenna}, reconstructed from the five
glued balls \texttt{I1..I5}. The fixture uses \emph{equal-length arms}, which makes the
optical constellation rotationally ambiguous; in this dataset the optical auto-labeling
swapped ball identities in ID01 and ID05.

Detection is rotation/translation invariant: for each capture, the ten pairwise distances
among \texttt{I1..I5} are compared against the clean 22-capture consensus fingerprint,
\[
[84.5,\,53.0,\,49.0,\,60.5,\,59.9,\,57.6,\,67.2,\,49.6,\,60.0,\,79.8]\mm,
\]
(pair order I1-I2,\,I1-I3,\,I1-I4,\,I1-I5,\,I2-I3,\,I2-I4,\,I2-I5,\,I3-I4,\,I3-I5,\,I4-I5),
with a clean-capture spread of $\approx\!3\mm$. A 120-permutation search uniquely identifies
the corrupted captures and their relabels:

\begin{table}[h]
\centering
\caption{Deterministic marker-label correction. Indices are the 0-based original ball
indices used for the new I1..I5.}
\begin{tabular}{lccc}
\toprule
ID & corrected I1..I5 & Motive$\to$corrected \texttt{Iantenna} shift & fingerprint dev.\ after relabel \\
\midrule
ID01 & $(0,1,4,2,3)$ & 54.1\mm & 1.5\mm \\
ID05 & $(3,4,2,0,1)$ & 2.1\mm  & 0.8\mm \\
\bottomrule
\end{tabular}
\end{table}

\noindent
All other 22 captures (including all eight low-position captures) are clean. ID04 is
explicitly clean despite a world-frame direction flip --- the correct diagnostic is the
ball-local fingerprint, not the world direction. Two consequences:
\begin{itemize}[leftmargin=1.4em]
  \item \textbf{ID01} ground truth was materially wrong (shift $54.1\mm$). After correction
        its UWB error does \emph{not} improve; it rises to $\approx\!139\mm$ in production
        \texttt{v4-io/all8}. The earlier (corrupted) truth had coincidentally sat closer to
        the AutoPos solution; the correction therefore \emph{reveals} a genuinely large
        error at ID01 rather than fixing it.
  \item \textbf{ID05} truth moved only $2.1\mm$ (the opposite-arm swap is symmetric about the
        antenna axis), so ID05's error is essentially real.
\end{itemize}
\textbf{Design lesson:} the next-generation fixture must use asymmetric, unequal arm lengths
so every ball has a unique distance signature in any pose; this ties into the planned
marker-hole PCB redesign.

\subsection{Production-output summary}
\begin{table}[h]
\centering
\caption{Production-output static-tag accuracy, \texttt{v4-io}.}
\begin{tabular}{lrrrrrr}
\toprule
eval set & median 3D & p95 3D & RMS 3D & median horiz.\ & median vert.\ & scale-bias diag.\ \\
\midrule
all8 & 77.4 & 270.3 & 138.3 & 43.8 & 63.1 & 30.0 \\
noG  & 81.3 & 278.6 & 141.1 & 46.9 & 63.5 & 31.2 \\
\bottomrule
\end{tabular}
\end{table}

\subsection{Raw replay matrix}
The raw replay matrix replays the raw static captures through the C-core T-series solvers
($5\;\text{anchor solvers} \times 4\;\text{tag solvers} \times \text{all8/noG}$).

\begin{table}[h]
\centering
\caption{Best all8 median 3D rows from the raw replay matrix.}
\begin{tabular}{rlllrrrr}
\toprule
rank & anchor & tag & eval & median 3D & p95 3D & RMS 3D & vert.\ median \\
\midrule
1 & v4-io   & T3 & all8 & 62.3 & 158.2 & 101.8 & 48.6 \\
2 & v4-io   & T4 & all8 & 69.1 & 182.3 & 107.0 & 55.0 \\
3 & v2      & T3 & all8 & 71.3 & 190.1 & 100.5 & 53.1 \\
4 & v3-lite & T3 & all8 & 72.4 & 190.2 & 100.6 & 52.7 \\
5 & v2      & T4 & all8 & 75.0 & 166.0 & 102.0 & 58.6 \\
\bottomrule
\end{tabular}
\end{table}

\begin{table}[h]
\centering
\caption{\texttt{v4-io/T4} all8 versus noG (raw replay).}
\begin{tabular}{lrrrrrr}
\toprule
eval set & median 3D & p95 3D & RMS 3D & horiz.\ median & vert.\ median & repeat.\ D3 median \\
\midrule
all8 & 69.1 & 182.3 & 107.0 & 41.3 & 55.0 & 67.4 \\
noG  & 83.9 & 291.1 & 138.7 & 67.9 & 43.9 & 56.2 \\
\bottomrule
\end{tabular}
\end{table}

\noindent
\textbf{Reading.} The best-case combination is \texttt{v4-io/T3/all8} (median $62.3\mm$); it
must be labelled as the best solver combination, not the representative production accuracy.
The deployment-oriented \texttt{v4-io/T4/all8} is close ($69.1\mm$). Dropping G does not
improve tag accuracy: for \texttt{v4-io/T4}, noG worsens median from $69.1$ to $83.9\mm$ and
p95 from $182.3$ to $291.1\mm$. This does not contradict the G-marker warning --- it means
that for tag solving the G \emph{UWB} measurement still helps more than the noG geometry loss
hurts, even though the G \emph{optical} marker is suspect.

\begin{figure}[h]
\centering
\begin{subfigure}{0.49\textwidth}\centering
  \includegraphics[width=\textwidth]{tag_error_by_position.png}
  \caption{Production tag error by position.}
\end{subfigure}\hfill
\begin{subfigure}{0.49\textwidth}\centering
  \includegraphics[width=\textwidth]{tag_error_vs_distance.png}
  \caption{Tag error versus distance.}
\end{subfigure}
\caption{Static-tag error by position and versus distance (production output).}
\label{fig:tagpos}
\end{figure}

\begin{figure}[h]
\centering
\begin{subfigure}{0.49\textwidth}\centering
  \includegraphics[width=\textwidth]{tag_raw_replay_accuracy_matrix.png}
  \caption{Raw replay accuracy matrix.}
\end{subfigure}\hfill
\begin{subfigure}{0.49\textwidth}\centering
  \includegraphics[width=\textwidth]{tag_raw_replay_v4io_by_position.png}
  \caption{Raw replay \texttt{v4-io} by position.}
\end{subfigure}
\caption{Raw replay matrix results.}
\label{fig:rawreplay}
\end{figure}

% ======================================================================
\section{Localization Metric Set (cm-scale)}
% ======================================================================
\label{sec:metrics}

\subsection{Which metrics apply, and which do not}
Standard localization checklists include real-time and dynamic metrics that do not apply to
this offline static dataset and are therefore \emph{deliberately omitted}: latency, update
rate, drop rate, availability/dropout (no real-time stream; offline replay of fixed
positions), and frame-to-frame jitter (a temporal/dynamic metric whose static analogue is
the already-reported repeatability D3 std). Crucially, 3D is reported as the \emph{primary}
metric and the vertical $Y$ axis is shown explicitly --- 2D is given only as a clearly
labelled secondary line. Demoting 3D to a supplement would hide the very axis (vertical) that
is the hard, headline part of pure-UWB 3D self-calibration.

\subsection{Full percentile set}
\begin{table}[h]
\centering
\caption{Full 3D-error percentile set (\si{\milli\meter}). Production \texttt{v4-io} and raw
replay \texttt{v4-io/T3} (best-case) and \texttt{T4} (deployment). Repeatability is the
static D3 std for cross-reference.}
\begin{tabular}{lrrrr}
\toprule
metric & prod.\ all8 & prod.\ noG & raw T3/all8 & raw T4/all8 \\
\midrule
mean        & 111.2 & 115.6 & 84.1  & 90.2  \\
RMSE        & 138.3 & 141.1 & 101.8 & 107.0 \\
P50         & 77.4  & 81.3  & 62.3  & 69.1  \\
P90         & 187.7 & 183.8 & 149.2 & 155.7 \\
P95         & 270.3 & 278.6 & 158.2 & 182.3 \\
P99         & 349.9 & 349.5 & 261.4 & 246.0 \\
max         & 369.6 & 365.7 & 292.2 & 263.8 \\
\midrule
repeat.\ D3 std P50 & 58.6 & 58.6 & 58.7  & 67.4 \\
repeat.\ D3 std P95 & 88.2 & 88.2 & 137.7 & 88.4 \\
\bottomrule
\end{tabular}
\end{table}

\noindent
The mean ($111\mm$) exceeds the median ($77\mm$), the classic ``mostly accurate, occasional
fliers'' signature.

\subsection{Per-axis signed bias --- the anisotropy result}
\begin{table}[h]
\centering
\caption{Per-axis signed bias and dispersion (OptiTrack frame, $Y$ vertical). Production
\texttt{v4-io/all8}. The error is strongly anisotropic: the vertical std is $\approx\!4\times$
the horizontal std, while all signed biases are small.}
\begin{tabular}{lrrr}
\toprule
component & signed mean bias & signed std & P95 abs.\ error \\
\midrule
$X$ (horizontal)        & $+23.7$ & 32.4  & 79.9  \\
$Y$ (\textbf{vertical}) & $-6.2$  & \textbf{130.8} & \textbf{259.4} \\
$Z$ (horizontal)        & $-8.8$  & 32.6  & 58.0  \\
\midrule
horizontal $XZ$ (2D, secondary) & \multicolumn{3}{l}{RMSE 51.6\quad P50 43.8\quad P95 82.8} \\
\bottomrule
\end{tabular}
\end{table}

\noindent
\textbf{This is the central error-structure result.} Horizontal P95 is only $58$--$80\mm$
(solid centimetre-level), whereas vertical P95 is $\approx\!260\mm$ with a std four times the
horizontal. Signed biases are all small ($|b|<25\mm$), so the error is \emph{not} a
systematic translation; it is a vertical-dominated dispersion. The raw-replay T3 solver
reduces the vertical std to $\approx\!82\mm$, indicating that the production solver is
specifically weaker in the vertical direction.

\subsection{Outlier rates (cm-scale thresholds)}
\begin{table}[h]
\centering
\caption{Fraction of the 24 positions within each 3D-error threshold. Thresholds are tuned
to this system's magnitude (50/80/100/200/300\,\si{\milli\meter}); generic metre-scale
thresholds are intentionally not used.}
\begin{tabular}{lrrrrr}
\toprule
within & prod.\ all8 & prod.\ noG & raw T3/all8 & raw T4/all8 \\
\midrule
$\le 50\mm$  & 25.0\% & 12.5\% & 33.3\% & 29.2\% \\
$\le 80\mm$  & 58.3\% & 50.0\% & 66.7\% & 62.5\% \\
$\le 100\mm$ & 58.3\% & 58.3\% & 70.8\% & 62.5\% \\
$\le 200\mm$ & 91.7\% & 91.7\% & 95.8\% & 95.8\% \\
$\le 300\mm$ & 95.8\% & 95.8\% & 100\%  & 100\%  \\
\bottomrule
\end{tabular}
\end{table}

\noindent
\textbf{Reading.} For production \texttt{v4-io/all8}, $58\%$ of positions are within
$80\mm$, $92\%$ within $200\mm$, and a single position exceeds $300\mm$. Notably the
$80\mm$ and $100\mm$ columns are identical (no position falls in $80$--$100\mm$): the
distribution is bimodal --- a cluster of good interior positions and a separate cluster of
edge/high positions pushed out by radial scale error --- rather than a single smooth spread.
Raw replay T3/T4 are tighter, with no position above $300\mm$.

\begin{figure}[h]
\centering
\begin{subfigure}{0.49\textwidth}\centering
  \includegraphics[width=\textwidth]{tag_error_cdf.png}
  \caption{3D position-error CDF.}
\end{subfigure}\hfill
\begin{subfigure}{0.49\textwidth}\centering
  \includegraphics[width=\textwidth]{tag_error_per_axis_bias.png}
  \caption{Per-axis signed bias and P95.}
\end{subfigure}
\caption{Localization metric visualizations: error CDF and per-axis bias.}
\label{fig:metrics}
\end{figure}

% ======================================================================
\section{Surveyed-Anchor Baseline: Isolating Self-Calibration Cost}
% ======================================================================
\label{sec:surveyed}

\subsection{Rationale and method}
The single AutoPos line cannot separate ``how much error comes from imperfect
self-calibration'' from ``the intrinsic UWB ranging/geometry floor.'' Most UWB papers have
no self-calibration step; they feed surveyed anchor coordinates into the tag solver. We add
exactly that as a control: OptiTrack-truth anchor coordinates are used directly as the
solver layout, the static tag is solved from the raw tag-to-anchor ranges, and the result is
compared directly to the corrected OptiTrack \texttt{Iantenna}. The entire path is in the
OptiTrack frame end to end --- \textbf{no Kabsch, no reflection, no scale fit} --- so it is a
clean control for the cost of replacing surveyed anchors with self-calibrated ones.

\subsection{Three delay treatments}
\begin{itemize}[leftmargin=1.4em]
  \item \texttt{raw\_zero\_delay}: raw ranges, all delays zero. Dominated by the missing
        endpoint delay; not the physical floor.
  \item \texttt{autopos\_v4io\_delay\_vector}: OptiTrack anchors with the \texttt{v4-io}
        AutoPos per-anchor delay vector. This is the hybrid that asks ``can the AutoPos delay
        be transplanted onto a surveyed layout?''.
  \item \texttt{inter\_anchor\_delaycal}: per-anchor endpoint delays fit from raw inter-anchor
        medians against OptiTrack true inter-anchor distances; tag delay set to the median
        endpoint delay. This uses OptiTrack \emph{twice} (anchors and delay), so it is partly
        circular and is an optimistic lower bound, not a field-achievable number.
\end{itemize}
The inter-anchor delaycal diagnostic has median common endpoint bias $84.3\mm$ and
per-anchor least-squares residual RMS $51.1\mm$.

\begin{table}[h]
\centering
\caption{Headline surveyed-anchor comparison against the production AutoPos line.}
\small
\begin{tabular}{llllrrrrr}
\toprule
eval & anchor source & delay & tag & median & p95 & RMS & h.med & v.med \\
\midrule
all8 & OptiTrack & raw\_zero      & T4 & 296.0 & 443.1 & 297.8 & 95.5 & 279.9 \\
all8 & OptiTrack & autopos\_v4io  & T4 & 241.9 & 376.3 & 240.4 & 79.9 & 214.8 \\
all8 & OptiTrack & delaycal       & T4 & 58.4  & 134.8 & 74.9  & 37.6 & 39.5 \\
all8 & AutoPos v4-io & estimated  & prod & 77.4 & 270.3 & 138.3 & 43.8 & 63.1 \\
\midrule
noG  & OptiTrack & raw\_zero      & T4 & 322.1 & 534.1 & 349.7 & 126.2 & 286.2 \\
noG  & OptiTrack & autopos\_v4io  & T4 & 255.1 & 471.0 & 287.7 & 114.5 & 206.0 \\
noG  & OptiTrack & delaycal       & T4 & 78.1  & 198.2 & 118.4 & 50.1  & 64.1 \\
noG  & AutoPos v4-io & estimated  & prod & 81.3 & 278.6 & 141.1 & 46.9 & 63.5 \\
\bottomrule
\end{tabular}
\end{table}

\subsection{The delay vector is coupled to the layout (by design)}
The hybrid row (OptiTrack layout + AutoPos delay, $241.9\mm$) is far worse than AutoPos used
as a whole ($77.4\mm$). This is \emph{not} a defect: AutoPos jointly estimates layout and
delay as a single self-consistent calibration, and the delay is absorbed into the system to
minimise ranging residuals against AutoPos's own geometry. Detaching the delay and pairing it
with a foreign (even perfect) layout breaks that self-consistency. The decomposition
(Section~\ref{sec:delay}) confirms the AutoPos delay's differential part is essentially
uncorrelated with the independently measured physical delay differences (Pearson
$r=-0.03$): the AutoPos delay is an \emph{effective} delay (true antenna delay plus geometry/
scale compensation), not a portable physical delay measurement. The correct framing is
positive: AutoPos performs a coupled layout-plus-delay self-calibration, which is exactly why
the integrated solution reaches $77\mm$ \emph{without any external delay measurement}.

\subsection{Worst-point resolution: the tail is layout, not UWB failure}
\begin{table}[h]
\centering
\caption{Worst-point resolution. The 270\,\si{\milli\meter}-class production tail collapses
under the OptiTrack-anchor delaycal control, locating it in the self-calibration/layout/
frame-lock budget rather than in irreducible UWB failure at those points.}
\begin{tabular}{lrrrl}
\toprule
ID & prod.\ all8 3D & surveyed+autopos-delay & surveyed delaycal & interpretation \\
\midrule
ID03 & 369.6 & 432.5 & 78.2 & collapses with OptiTrack delaycal \\
ID04 & 283.9 & 388.1 & 89.5 & collapses with OptiTrack delaycal \\
ID06 & 192.8 & 244.5 & 34.6 & collapses with OptiTrack delaycal \\
ID05 & 52.1  & 145.2 & 32.6 & moderate; truth shift was only 2.1\,\si{\milli\meter} \\
\bottomrule
\end{tabular}
\end{table}

\noindent
\textbf{Important caveat.} The delaycal lower bound ($58\mm$) uses OptiTrack truth twice and
is \emph{not} a number AutoPos can reach independently. The defensible conclusion is about
\emph{error attribution}: the large tail arises in the calibration layer (anchor coordinates
and delay), not in the UWB ranging at those points. It is \emph{not} that ``AutoPos is only
$18.9\mm$ from optimal'' --- that gap is measured against an OptiTrack-dependent floor.

\begin{figure}[h]
\centering
\begin{subfigure}{0.49\textwidth}\centering
  \includegraphics[width=\textwidth]{autopos_vs_surveyed_per_position.png}
  \caption{AutoPos versus surveyed, per position.}
\end{subfigure}\hfill
\begin{subfigure}{0.49\textwidth}\centering
  \includegraphics[width=\textwidth]{worst_points_autopos_vs_surveyed.png}
  \caption{Worst points, AutoPos versus surveyed.}
\end{subfigure}
\caption{Surveyed-anchor control.}
\label{fig:surveyed}
\end{figure}

% ======================================================================
\section{VDOP Geometry}
% ======================================================================
\label{sec:vdop}

The default geometric model is the range-only Jacobian row $[u_x,u_y,u_z]$, matching the tag
solver (which does not co-estimate a clock/range bias). The range-bias model
$[u_x,u_y,u_z,1]$ is computed only as a diagnostic appendix and is not the default
explanation, because adding the bias column inflates VDOP and does not match the solver.

\begin{table}[h]
\centering
\caption{Range-only VDOP summary on the 25\,\si{\milli\meter} grid.}
\begin{tabular}{lrrrrr}
\toprule
mask & VDOP median & VDOP p90 & VDOP p95 & GDOP median & HDOP median \\
\midrule
all8  & 0.804 & 0.911 & 0.941 & 1.170 & 0.861 \\
noG   & 0.856 & 1.057 & 1.177 & 1.278 & 0.943 \\
dropH & 0.861 & 1.075 & 1.199 & 1.283 & 0.941 \\
\bottomrule
\end{tabular}
\end{table}

\noindent
\textbf{Reading.} The all8 geometry is healthy across the working volume. Dropping an upper
anchor (G or H) raises VDOP and condition number but does not cause rank failure. Critically,
the worst-error tag positions (ID03/ID04/ID06) do \emph{not} coincide with the worst VDOP ---
in fact they sit at favourable (low) VDOP. Therefore geometry alone does not explain the tail;
combined with Sections~\ref{sec:tag} and~\ref{sec:surveyed}, the tail is attributed to
layout/scale error, not geometric dilution.

\begin{figure}[h]
\centering
\begin{subfigure}{0.49\textwidth}\centering
  \includegraphics[width=\textwidth]{vdop_slices_grid25.png}
  \caption{VDOP slices (25\,\si{\milli\meter} grid).}
\end{subfigure}\hfill
\begin{subfigure}{0.49\textwidth}\centering
  \includegraphics[width=\textwidth]{dop_4panel_mid_all8_grid25.png}
  \caption{GDOP/PDOP/HDOP/VDOP at mid height.}
\end{subfigure}
\caption{Range-only DOP geometry.}
\label{fig:vdop}
\end{figure}

% ======================================================================
\section{Robustness: MC5000 Random Keep-$k$}
% ======================================================================
\label{sec:mc}

The completed matrix is $5\;\text{anchor solvers}\times 4\;\text{tag solvers}\times
\text{static/roto}\times\text{keep }8/7/6/5/4\times\text{MC5000}$; integrity is 40/40 blocks
PASS with no issues.

\begin{table}[h]
\centering
\caption{\texttt{v4-io/T4} random keep-$k$ degradation (medians).}
\begin{tabular}{lrrrrr}
\toprule
kind & keep8 & keep7 & keep6 & keep5 & keep4 \\
\midrule
static D3 std median (\si{\milli\meter})    & 61.7 & 83.2 & 107.3 & 149.1 & 196.4 \\
roto turn-center RMS median (\si{\milli\meter}) & 12.1 & 21.6 & 30.5  & 41.4  & 63.9  \\
\bottomrule
\end{tabular}
\end{table}

\noindent
Degradation is monotonic and physically plausible. Note that the random keep-7 \emph{median}
(averaging over all single-drop choices) is not the same as any specific single-anchor drop;
the latter is treated in Section~\ref{sec:anchorhealth} and can individually be better or
worse than keep-8 depending on which anchor is removed.

\begin{figure}[h]
\centering
\begin{subfigure}{0.49\textwidth}\centering
  \includegraphics[width=\textwidth]{mc_keepk_static_curves.png}
  \caption{Static keep-$k$ curves.}
\end{subfigure}\hfill
\begin{subfigure}{0.49\textwidth}\centering
  \includegraphics[width=\textwidth]{mc_keepk_roto_curves.png}
  \caption{Roto keep-$k$ curves.}
\end{subfigure}
\caption{MC5000 keep-$k$ robustness.}
\label{fig:mc}
\end{figure}

% ======================================================================
\section{Robustness: Stratified Fixed-Drop Keep-$k$}
% ======================================================================
\label{sec:stratified}

This exhaustive fixed dropped-set replay (6480 detail rows; 40/40 blocks; 162 fixed
keep-sets per block) answers ``which specific anchors are dangerous to lose,'' which the
random MC cannot.

\begin{table}[h]
\centering
\caption{\texttt{v4-io/T4} composition snapshot (lower-heavy vs upper-heavy drops).}
\begin{tabular}{lrrrl}
\toprule
kind & keep$k$ & lower-heavy & upper-heavy & interpretation \\
\midrule
static D3 std    & 7 & 65.2 & 61.0 & lower drop slightly worse \\
static D3 std    & 6 & 84.5 & 63.1 & lower drop worse \\
static D3 std    & 5 & 96.0 & 65.1 & lower drop worse \\
static D3 std    & 4 & 117.3 & 57.2 & lower drop much worse \\
roto center RMS  & 7 & 18.3 & 16.0 & lower drop slightly worse \\
roto center RMS  & 6 & 25.2 & 18.1 & lower drop worse \\
roto center RMS  & 5 & 30.9 & 24.1 & lower drop worse \\
roto center RMS  & 4 & 48.5 & 34.0 & lower drop worse \\
\bottomrule
\end{tabular}
\end{table}

\noindent
\textbf{Reading.} The naive intuition ``dropping upper anchors hurts $Z$ most'' is \emph{not}
what the data shows. In this geometry, lower-heavy drops are more damaging for both static
repeatability and roto stability. Balanced drops can still be catastrophic (the worst roto
keep-4 set is \texttt{ACEG}), so anchor identity matters, not only the remaining count.

\begin{figure}[h]
\centering
\includegraphics[width=0.7\textwidth]{stratified_keepk_upper_vs_lower.png}
\caption{Stratified keep-$k$: upper-heavy versus lower-heavy drops.}
\label{fig:stratified}
\end{figure}

% ======================================================================
\section{Per-Anchor Diagnostics}
% ======================================================================
\label{sec:pair}

\subsection{Pair residuals and raw asymmetry}
\begin{table}[h]
\centering
\caption{\texttt{v4-io} worst all-1000 residual pairs and largest raw directional
asymmetries.}
\begin{tabular}{lrl@{\hspace{2em}}lr}
\toprule
pair & residual (\si{\milli\meter}) & involves G & pair & asymmetry (\si{\milli\meter}) \\
\midrule
B-C & $-141.7$ & no  & D-F & 22.0 \\
B-G & $+113.7$ & yes & E-F & 20.0 \\
D-E & $+99.8$  & no  & C-F & 18.0 \\
D-F & $+91.4$  & no  & A-H & $-17.0$ \\
F-H & $-59.3$  & no  & B-E & $-16.0$ \\
\bottomrule
\end{tabular}
\end{table}

\noindent
The worst residual pairs are not only G-involving; B-C, D-E, D-F and F-H show that the
residual structure is broader than a single-anchor G issue. Directional asymmetry exists but
is modest relative to the largest layout residuals.

\begin{figure}[h]
\centering
\begin{subfigure}{0.32\textwidth}\centering
  \includegraphics[width=\textwidth]{pair_residual_bias_heatmap.png}
  \caption{Residual bias.}
\end{subfigure}\hfill
\begin{subfigure}{0.32\textwidth}\centering
  \includegraphics[width=\textwidth]{pair_residual_abs_heatmap.png}
  \caption{Residual magnitude.}
\end{subfigure}\hfill
\begin{subfigure}{0.32\textwidth}\centering
  \includegraphics[width=\textwidth]{pair_raw_asymmetry_heatmap.png}
  \caption{Raw asymmetry.}
\end{subfigure}
\caption{Pair-residual and raw-sweep diagnostics.}
\label{fig:pairheat}
\end{figure}

\subsection{Temporal/thermal drift}
A raw-link drift diagnostic over 24 static sessions (192 anchor-session links): median
absolute drift slope $1.54\mm/\text{min}$, p95 $16.21\mm/\text{min}$; median absolute drift
over a capture $3.07\mm$, p95 $32.42\mm$. Per-anchor, A--E drift only a few mm over a
capture, while F/G/H have larger tails (G and F the most visible; G p95 slope
$\approx\!24.7\mm/\text{min}$, F $\approx\!19.2\mm/\text{min}$). Drift is not the dominant
explanation for the median absolute error, but F/G/H are placement/hardware-sensitive and
relevant to the height-alignment story.

\begin{figure}[h]
\centering
\begin{subfigure}{0.32\textwidth}\centering
  \includegraphics[width=\textwidth]{temporal_drift_slope_heatmap.png}
  \caption{Drift slope heatmap.}
\end{subfigure}\hfill
\begin{subfigure}{0.32\textwidth}\centering
  \includegraphics[width=\textwidth]{temporal_drift_slope_boxplot.png}
  \caption{Drift slope boxplot.}
\end{subfigure}\hfill
\begin{subfigure}{0.32\textwidth}\centering
  \includegraphics[width=\textwidth]{temporal_drift_worst_timeseries.png}
  \caption{Worst-link time series.}
\end{subfigure}
\caption{Temporal/thermal drift diagnostics.}
\label{fig:drift}
\end{figure}

\subsection{Anchor health scorecard and single-anchor criticality}
\label{sec:anchorhealth}
A heuristic composite trust score combines pair-residual magnitude, raw asymmetry, temporal
drift, OptiTrack marker-fingerprint flag, and delay differential. Least trusted: G (0.46), D
(0.56), H (0.60), F (0.62); most trusted: B (0.86), E (0.77), A (0.71). The single-anchor
criticality (drop-one) ranking is dominated by noise: all single-drop deltas are within
$\pm11\mm$, which is smaller than the repeatability CI, so no anchor shows a statistically
distinguishable dominant role --- consistent with the near-symmetric box layout and the
all-faces test design. The one weakly real direction is that dropping the high-drift anchors
F/G tends to slightly \emph{improve} static repeatability (they add ranging noise while
contributing geometric constraint), mirroring the all8-vs-noG behaviour in
Section~\ref{sec:tag}.

\begin{figure}[h]
\centering
\begin{subfigure}{0.49\textwidth}\centering
  \includegraphics[width=\textwidth]{anchor_health_scorecard.png}
  \caption{Anchor health scorecard.}
\end{subfigure}\hfill
\begin{subfigure}{0.49\textwidth}\centering
  \includegraphics[width=\textwidth]{single_anchor_criticality.png}
  \caption{Single-anchor criticality.}
\end{subfigure}
\caption{Per-anchor health and criticality. Criticality differences are within measurement
uncertainty.}
\label{fig:anchorhealth}
\end{figure}

% ======================================================================
\section{Additional Error-Structure Diagnostics}
% ======================================================================
\label{sec:errstruct}

\subsection{Antenna-delay common/differential decomposition}
\label{sec:delay}
Decomposing the \texttt{v4-io} delay vector $d_i = \overline{d} + (d_i-\overline{d})$ gives a
common term $\overline{d}\approx34\mm$ and a differential spread (std $\approx\!19\mm$, range
$-34$ to $+26\mm$). Compared against the independent inter-anchor delaycal endpoint delays,
the \emph{differential} parts are essentially uncorrelated (Pearson $r=-0.03$, RMSE
$\approx\!38\mm$), while the common terms differ ($34$ vs $91\mm$). This confirms that the
AutoPos delay is an effective joint-self-calibration parameter (true antenna delay plus
geometry/scale compensation), coupled to the layout gauge by design, not a standalone
physical antenna-delay measurement.

\subsection{Scale propagation (error versus distance from centre)}
Regressing tag error on distance from the array centroid gives, for the radial component, a
slope $\approx\!230\mm/\text{m}$ with $R^2=0.48$ ($p=2\times10^{-4}$); the 3D error gives
$R^2=0.28$ ($p=0.007$). Both are positive and significant: error grows with distance from
centre, the signature of the $\approx\!4\%$ scale bias.

\subsection{Error vector field (direction, not magnitude)}
The mean error vector is small ($\approx\!26\mm$ norm) relative to the scatter
($\approx\!136\mm$), giving a mean-to-scatter ratio of $0.19$; the median radial component
($63\mm$) exceeds the median tangential ($40\mm$); and $83\%$ of error vectors point radially
\emph{outward}. Together: the error is predominantly radial-outward (scale), not a common
translation (small mean vector) and not isotropic noise (directional preference).

\subsection{Worst-point raw-range fingerprint}
For ID01/ID03/ID04/ID06 the per-anchor raw-range residuals are large but not concentrated on
a single anchor, and clean reference positions (e.g.\ ID22) also show large single-anchor
residuals while remaining accurate overall. This rules out ``one bad anchor'' and is
consistent with geometric amplification of scale error at edge/high positions.

\subsection{Stratified tag error (height, edge/centre, facing)}
These small-group splits ($n$ of a few per group, so exploratory only) are consistent with
the scale-propagation picture: edge and high positions are worse, centre/mid better, and the
facing groups show no strong VDOP-linked differentiation.

\begin{figure}[h]
\centering
\begin{subfigure}{0.32\textwidth}\centering
  \includegraphics[width=\textwidth]{delay_decomposition.png}
  \caption{Delay decomposition.}
\end{subfigure}\hfill
\begin{subfigure}{0.32\textwidth}\centering
  \includegraphics[width=\textwidth]{tag_error_vs_center_distance.png}
  \caption{Error vs distance from centre.}
\end{subfigure}\hfill
\begin{subfigure}{0.32\textwidth}\centering
  \includegraphics[width=\textwidth]{tag_error_vector_field.png}
  \caption{Error vector field.}
\end{subfigure}
\caption{Error-structure diagnostics: delay decomposition, scale propagation, vector field.}
\label{fig:errstruct}
\end{figure}

\begin{figure}[h]
\centering
\begin{subfigure}{0.32\textwidth}\centering
  \includegraphics[width=\textwidth]{tag_error_by_height.png}
  \caption{By height.}
\end{subfigure}\hfill
\begin{subfigure}{0.32\textwidth}\centering
  \includegraphics[width=\textwidth]{tag_error_edge_vs_center.png}
  \caption{Edge vs centre.}
\end{subfigure}\hfill
\begin{subfigure}{0.32\textwidth}\centering
  \includegraphics[width=\textwidth]{tag_error_by_facing.png}
  \caption{By facing group.}
\end{subfigure}
\caption{Stratified tag error (exploratory; small group sizes).}
\label{fig:straterr}
\end{figure}

\begin{figure}[h]
\centering
\includegraphics[width=0.75\textwidth]{worstpoint_range_residual_fingerprint.png}
\caption{Worst-point raw-range residual fingerprint versus clean references.}
\label{fig:worstfinger}
\end{figure}

\subsection{Synthesis: what dominates the tag error}
The combined evidence identifies a single dominant cause:
\begin{enumerate}[leftmargin=1.6em]
  \item errors point radially outward ($83\%$) and grow with distance from centre
        ($R^2=0.48$, $p=2\times10^{-4}$) --- a scale signature;
  \item there is no common translation (mean-vector/scatter ratio $0.19$);
  \item the AutoPos delay differential is uncorrelated with physical delay ($r=-0.03$), i.e.\
        scale is absorbed into the effective delay (layout--delay coupling);
  \item worst-point raw residuals are not concentrated on one anchor, ruling out a single bad
        anchor;
  \item the surveyed-anchor control shows the tail collapses with perfect anchors, locating
        it in the calibration layer rather than the UWB ranging floor.
\end{enumerate}
\textbf{Conclusion:} the tag error is dominated by the self-calibrated layout's
$\approx\!4\%$ scale bias propagating radially, with the scale coupled to the antenna-delay
estimate; it is \emph{not} dominated by a single weak anchor or by isotropic measurement
noise. The clear improvement path is to break the scale/delay coupling with an independent
baseline (tape/laser) delay calibration, rather than to change the UWB ranging or swap an
anchor.

\begin{figure}[h]
\centering
\includegraphics[width=0.85\textwidth]{bootstrap_confidence_intervals.png}
\caption{Bootstrap 95\% confidence intervals for headline metrics.}
\label{fig:bootstrap}
\end{figure}

\noindent
\textbf{Confidence intervals.} Bootstrap ($n_\text{boot}=10000$, seed 42) gives wide CIs for
the 24-position static summaries (e.g.\ tag absolute median $77.4\mm$, CI $52.1$--$141.1$;
p95 $270.3\mm$, CI $151.9$--$369.6$; layout rigid RMS $104.9\mm$, CI $81.4$--$133.2$). Median
differences of a few mm should not be over-interpreted; MC repeat intervals are tighter
because they derive from the MC5000 repeat distribution.

% ======================================================================
\section{Rotational (Roto) Absolute Validation}
% ======================================================================
\label{sec:roto}

\todo{RESERVED SECTION --- TO BE COMPLETED}

\noindent
This section is intentionally reserved so that the rotational validation is not forgotten.
At the time of writing, only UWB-only consistency diagnostics are available; the independent
OptiTrack absolute validation of the rotational captures is pending external processing.

\subsection{Current UWB-only consistency (not absolute accuracy)}
\begin{table}[h]
\centering
\caption{\texttt{v4-io} roto UWB-only consistency (bootstrap 95\% CI). These are consistency
diagnostics, NOT OptiTrack-validated absolute accuracy.}
\begin{tabular}{lrrr}
\toprule
metric & point (\si{\milli\meter}) & CI low & CI high \\
\midrule
roto abs $\Delta R$ error median   & 33.33 & 22.52 & 40.15 \\
roto abs $\Delta R$ error p95      & 42.58 & 40.17 & 47.25 \\
roto turn-center RMS median        & 14.31 & 13.16 & 17.36 \\
\bottomrule
\end{tabular}
\end{table}

\noindent
The turn-center RMS ($14.31\mm$) is tight, but it mostly tests shape/periodic consistency
about a fitted trajectory, not absolute alignment to OptiTrack. \textbf{No absolute roto
accuracy may be claimed until the OptiTrack roto ground truth is processed.}

\subsection{To be completed when OptiTrack roto data lands}
\begin{itemize}[leftmargin=1.4em]
  \item Absolute roto trajectory error against OptiTrack (ATE/RMSE, per-axis, CDF).
  \item Turn-center / axis-of-rotation absolute accuracy versus OptiTrack.
  \item Dynamic-versus-static comparison and any motion-induced degradation.
  \item Roto-specific outlier and latency behaviour if a real-time stream is recorded.
  \item \todo{Insert roto figures here once generated.}
\end{itemize}

% ======================================================================
\section{Limitations and Open Items}
% ======================================================================
\label{sec:limits}

\begin{itemize}[leftmargin=1.4em]
  \item \textbf{Roto absolute validation pending} (Section~\ref{sec:roto}); only UWB-only
        consistency is available.
  \item \textbf{Delay/scale circularity.} OptiTrack-derived delay calibration and
        OptiTrack-based layout validation are not independent if used as a closed loop. The
        clean path is an independent tape/laser baseline delay calibration, frozen before
        OptiTrack validation. Static-tag validation remains independent as long as the tag
        truth is not used for delay calibration, scale fitting, or solver selection.
  \item \textbf{Anchor G optical marker} is suspect (fingerprint), but changes the layout RMS
        by only $\approx\!0.5\mm$; all8/noG dual reporting is retained for rigor, not because
        G changes the conclusion.
  \item \textbf{Marker fixture} used equal-length arms (caused the ID01/ID05 swaps); the
        next fixture must use asymmetric arm lengths.
  \item \textbf{Static-tag p95} remains wide and placement-sensitive; the tail is a
        layout/scale effect concentrated at edge/high positions.
  \item \textbf{Small-group analyses} (height/edge/facing) are exploratory due to few points
        per group.
\end{itemize}

% ======================================================================
\section{Recommended Headline Statements}
% ======================================================================
\label{sec:headline}

\paragraph{Anchor layout.}
Using reflection-allowed rigid alignment with no scale fit, the \texttt{v4-io} anchor layout
reaches $104.9\mm$ 3D RMS against OptiTrack across all eight anchors (horizontal $86.1\mm$,
vertical $59.9\mm$); the similarity scale diagnostic is $0.960$ and is not an accuracy claim.

\paragraph{Static tag.}
Using the anchor-locked transform and corrected ground truth, production \texttt{v4-io/all8}
is $77.4\mm$ median 3D and $270.3\mm$ p95; the error is strongly anisotropic (horizontal p95
$\approx\!58$--$80\mm$, vertical p95 $\approx\!260\mm$) with no significant signed bias. The
best-case raw-replay combination is \texttt{v4-io/T3/all8} ($62.3\mm$ median, $158.2\mm$ p95);
the deployment-oriented \texttt{v4-io/T4/all8} is $69.1\mm$ median, $182.3\mm$ p95. $58\%$ of
positions are within $80\mm$ and $92\%$ within $200\mm$.

\paragraph{Surveyed-anchor control.}
With OptiTrack anchors and an OptiTrack-derived inter-anchor delay calibration, the all8/T4
lower bound is $58.4\mm$ median; the production tail (ID03/ID04/ID06) collapses under this
control, attributing the tail to self-calibration/layout rather than UWB failure. The AutoPos
delay vector does not transfer to a surveyed layout ($241.9\mm$), confirming a coupled,
by-design joint self-calibration.

\paragraph{Robustness.}
MC5000 completed 40/40 blocks. For \texttt{v4-io/T4}, static D3 std median degrades
$61.7\to196.4\mm$ from keep8 to keep4. Stratified analysis shows lower-heavy drops are more
damaging than upper-heavy in this geometry; single-anchor criticality differences are within
measurement uncertainty.

\paragraph{Roto.}
UWB-only turn-center RMS median is $14.31\mm$ (CI $13.16$--$17.36$), but roto absolute
OptiTrack validation is pending and no absolute roto accuracy is claimed.

\paragraph{Positioning of the contribution.}
The strength of this work is not the smallest number but the hardest problem definition (pure
3D inter-anchor self-calibration with no external sensor aid), the strongest independent
validation (OptiTrack), and the most decomposed error budget. The defensible claim is
credible decimetre-level 3D anchor calibration and approximately $6$--$8\,\si{\centi\meter}$
static absolute capture inside the working volume, not a general centimetre-level
optical-mocap replacement.

\end{document}
